\section{Related Work}
\label{sec:related}

\paragraph{LLM benchmarks.}
General-purpose benchmarks such as MMLU \citep{hendrycks2021mmlu},
HellaSwag \citep{zellers2019hellaswag}, and BIG-Bench
\citep{srivastava2023bigbench} report aggregate accuracy over large,
static item sets. They are invaluable for coarse ranking but do not
provide per-item difficulty parameters, standard errors, or a latent
scale on which to compare interventions. \todo{Add 1--2 domain-specific
medical/regulatory benchmarks for contrast.}

\paragraph{IRT for evaluating machines.}
A growing line of work applies item-response theory to model evaluation,
using IRT to identify informative items, to weight benchmark scores, and
to compare systems on a latent ability scale
\citep[e.g.][]{lalor2019pyirt}. \todo{Cite the specific
``IRT-for-NLP-evaluation'' papers you want to position against; note
whether they reuse human exams vs. build machine-native instruments.}
Our work differs in that the items are \emph{generated and calibrated for
machine examinees} in a domain with no human calibration pool, and the
difficulty mandate is built into generation.

\paragraph{IRT in educational measurement.}
The Rasch (1PL) model \citep{rasch1960} and the broader IRT family
\citep{lord1980} are the foundation of modern high-stakes testing.
We adopt Rasch for its specific objectivity---ability differences are
invariant to which items are used---which is precisely the property we
need when comparing $\theta$ across prompt interventions
(Section~\ref{sec:deltatheta}). We use the \texttt{girth} library
\citep{girth} for marginal-maximum-likelihood estimation and
\texttt{py-irt} \citep{lalor2019pyirt} for posterior uncertainty.

\paragraph{LLM-as-judge.}
Using a strong LLM to grade open-ended responses is now common practice
\citep{zheng2023judge}. We adopt it out of necessity---there is no human
grading pool---and mitigate its known biases with a domain-specific
rubric and a strict/lenient dual-judge design
(Section~\ref{sec:itemgen}). We treat judge error as a measured quantity
rather than assuming it away (Section~\ref{sec:discussion}).
